% Report on the second experiment, 23 November 2026.
%
%   .venv/bin/python v2/writeup/figures.py                       (figures and numbers)
%   cd v2/writeup && ~/.local/share/tectonic-0.17.0/tectonic report.tex
%
% Every number in the text is a macro from numbers.tex, which figures.py writes from the run
% files. Nothing is typed by hand, so text and figures cannot disagree.

\documentclass[10pt,letterpaper]{article}
\usepackage[letterpaper,textwidth=6.4in,top=0.68in,bottom=0.86in,footskip=0.4in]{geometry}

% One family of typefaces for text and figures: IBM Plex. Serif for the text, Sans for headings,
% captions, tables and everything inside the figures, Mono for words of the made-up language.
\usepackage{fontspec}
\setmainfont{IBMPlexSerif}[Path=fonts/,Extension=.otf,UprightFont=*-Regular,ItalicFont=*-Italic,
  BoldFont=*-SemiBold,BoldItalicFont=*-BoldItalic]
\setsansfont{IBMPlexSans}[Path=fonts/,Extension=.otf,UprightFont=*-Regular,ItalicFont=*-Italic,
  BoldFont=*-SemiBold]
\setmonofont{IBMPlexMono}[Path=fonts/,Extension=.otf,UprightFont=*-Regular,BoldFont=*-SemiBold,Scale=0.9]
\linespread{1.09}

\usepackage{microtype}
\usepackage{graphicx}
\usepackage{float}
\usepackage{flafter}   % a figure never appears before the place where the text brings it up
\usepackage{needspace}
\usepackage{xcolor}
\usepackage{booktabs,array}
\usepackage{enumitem}
\usepackage{titlesec}
\usepackage[most]{tcolorbox}
\usepackage{fancyhdr}
\usepackage[font={sf,footnotesize},labelfont=bf,labelsep=period,skip=5pt,justification=raggedright,
  singlelinecheck=false]{caption}
\usepackage[hidelinks,pdftitle={Does a larger vocabulary need a larger model? Second experiment}]{hyperref}

% The colours of the figures (see the top of figures.py).
\definecolor{ink}{HTML}{1A1A19}
\definecolor{second}{HTML}{52514E}
\definecolor{muted}{HTML}{898781}
\definecolor{hair}{HTML}{C3C2B7}
\definecolor{pale}{HTML}{F4F3EF}
\color{ink}

\setlength{\parindent}{0pt}
\setlength{\parskip}{4.6pt plus 1pt minus 0.5pt}
\setlist{itemsep=2pt,topsep=1.5pt,parsep=0pt,partopsep=0pt,leftmargin=1.25em}
\raggedbottom
\widowpenalty=10000 \clubpenalty=10000
\emergencystretch=1.5em

% Headings. A finding is a full sentence under a small label; so is the part on the predictions,
% which stands beside the findings.
\titleformat{\section}{\sffamily\bfseries\fontsize{14.5}{18}\selectfont}{\textcolor{muted}{\thesection}}{0.55em}{}
\titlespacing*{\section}{0pt}{17pt plus 3pt minus 2pt}{3pt}
\titleformat{\subsection}{\sffamily\bfseries\fontsize{10.5}{13}\selectfont}{}{0pt}{}
\titlespacing*{\subsection}{0pt}{9pt plus 2pt minus 1pt}{0pt}
\newcommand{\smallcaps}[1]{{\sffamily\fontsize{7.6}{9}\selectfont\color{second}\addfontfeature{LetterSpace=7}#1}}
\newcommand{\labelled}[2]{%
  \vspace{12pt plus 3pt minus 2pt}%
  \smallcaps{#1}\par\nopagebreak\vspace{-2pt}%
  {\sffamily\bfseries\fontsize{12.5}{15.5}\selectfont\raggedright #2\par}%
  \nopagebreak\vspace{1pt}}
\newcounter{finding}
\newcommand{\finding}[1]{\refstepcounter{finding}\labelled{FINDING \thefinding}{#1}}

\newcommand{\s}[1]{\texttt{#1}}                  % a word or sentence of the made-up language
\newcommand{\term}[1]{\textbf{\textit{#1}}}      % a term at the place where it is explained
\newcommand{\lead}[1]{\textbf{#1}}               % the opening words of a point
\newcommand{\point}[1]{{\color{ink}\bfseries #1}} % the first sentence of a caption: the point
\newcommand{\see}[1]{\mbox{\sffamily\footnotesize\color{second}(#1)}} % where an answer on page one is shown
\captionsetup{textfont={color=second}}
\setlength{\intextsep}{9pt plus 2pt minus 1pt}
\setlength{\textfloatsep}{11pt plus 2pt minus 1pt}
\renewcommand{\topfraction}{0.9}\renewcommand{\bottomfraction}{0.9}\renewcommand{\textfraction}{0.05}
\renewcommand{\floatpagefraction}{0.85}\setcounter{topnumber}{2}\setcounter{totalnumber}{4}
% A figure at full width with its caption underneath. \fig puts it where the text brings it up if
% it fits on the page, and otherwise at the top of the next page. \fighere puts it exactly where
% it is called: the findings are laid out so that each of their figures closes a page, with the
% text that reads it just before and just after.
\newcommand{\fig}[3]{\begin{figure}[!ht]\includegraphics[width=\linewidth]{figures/#1.pdf}%
  \caption{#3}\label{#2}\end{figure}}
\newcommand{\fighere}[3]{\begin{figure}[H]\includegraphics[width=\linewidth]{figures/#1.pdf}%
  \caption{#3}\label{#2}\end{figure}}
\newcommand{\none}{\textcolor{muted}{--}}
\newcommand{\reached}[1]{\textbf{#1}}

\pagestyle{fancy}\fancyhf{}\renewcommand{\headrulewidth}{0pt}
\fancyfoot[L]{\sffamily\fontsize{7.6}{9}\selectfont\color{muted}Does a larger vocabulary need a larger model?
  \enspace·\enspace Second experiment \enspace·\enspace 23 November 2026}
\fancyfoot[R]{\sffamily\fontsize{7.6}{9}\selectfont\color{second}\thepage}

\newcommand{\CloudCost}{24.39}
\newcommand{\CloudMachines}{100}
\newcommand{\RunsOnOwnComputer}{48}
\newcommand{\PracticeEight}{15}
\newcommand{\PracticeSixteen}{4}
\newcommand{\PracticeThirtyTwo}{4}
\newcommand{\PracticeSixtyFour}{6}
\newcommand{\TimingRuns}{Four}
\newcommand{\PlannedSixtyFour}{120}
\newcommand{\SentencesPerStep}{512}
\newcommand{\ThirdRuns}{2,880}
\newcommand{\LongerStartRunsIn}{34}
\newcommand{\KnownBar}{95}
\newcommand{\AxisStartSettings}{70}
\newcommand{\AxisStartCurves}{70}
\newcommand{\KeyBarsWord}{eight}
\newcommand{\PracticeRuns}{29}
\newcommand{\QuestionExampleRight}{83}
\newcommand{\QuestionExampleScore}{66}
\newcommand{\RunsCounted}{476}
\newcommand{\RunsLonger}{8}
\newcommand{\RunsLongerWord}{eight}
\newcommand{\RunsAll}{484}
\newcommand{\StagesWord}{four}
\newcommand{\CreatureTimes}{64}
\newcommand{\RunsEight}{200}
\newcommand{\RepeatsEight}{5}
\newcommand{\RepeatsEightWord}{five}
\newcommand{\StepsEight}{40,000}
\newcommand{\RateEight}{0.01}
\newcommand{\SentencesEight}{20.5}
\newcommand{\RunsSixteen}{120}
\newcommand{\RepeatsSixteen}{3}
\newcommand{\RepeatsSixteenWord}{three}
\newcommand{\StepsSixteen}{40,000}
\newcommand{\RateSixteen}{0.01}
\newcommand{\SentencesSixteen}{20.5}
\newcommand{\RunsThirtyTwo}{120}
\newcommand{\RepeatsThirtyTwo}{3}
\newcommand{\RepeatsThirtyTwoWord}{three}
\newcommand{\StepsThirtyTwo}{60,000}
\newcommand{\RateThirtyTwo}{0.01}
\newcommand{\SentencesThirtyTwo}{30.7}
\newcommand{\RunsSixtyFour}{36}
\newcommand{\RepeatsSixtyFour}{2}
\newcommand{\RepeatsSixtyFourWord}{two}
\newcommand{\StepsSixtyFour}{160,000}
\newcommand{\RateSixtyFour}{0.003}
\newcommand{\SentencesSixtyFour}{81.9}
\newcommand{\RunsEightSixteen}{320}
\newcommand{\RepeatsFewest}{2}
\newcommand{\RepeatsMost}{5}
\newcommand{\StepsLonger}{640,000}
\newcommand{\SentencesLonger}{327.7}
\newcommand{\LongerTimesWord}{four}
\newcommand{\LongerTimesWordCap}{Four}
\newcommand{\LongerHoursLow}{6.6}
\newcommand{\LongerHoursHigh}{12.6}
\newcommand{\StepsShared}{144,815}
\newcommand{\StepsLate}{102,400}
\newcommand{\FirstSteps}{100, 141, 200, 283}
\newcommand{\StepRatio}{1.414}
\newcommand{\StepRatioShort}{1.4}
\newcommand{\CountedLastSteps}{102,400, 144,815 and 160,000}
\newcommand{\LongerLastSteps}{144,815, 204,800, 289,631, 409,600, 579,262 and 640,000}
\newcommand{\RowsSmall}{47}
\newcommand{\RowsLarge}{1,055}
\newcommand{\TableRowsEight}{1,055}
\newcommand{\TableRowsSixteen}{1,075}
\newcommand{\TableRowsThirtyTwo}{1,115}
\newcommand{\TableRowsSixtyFour}{1,195}
\newcommand{\RowsNeverChangedWordCap}{Four}
\newcommand{\OtherWords}{30}
\newcommand{\WordsPerKind}{five}
\newcommand{\OutsideNarrow}{11,278}
\newcommand{\OutsideWide}{26,878}
\newcommand{\TableNumbersLarge}{33,760}
\newcommand{\ModelNarrow}{45,038}
\newcommand{\ModelWide}{60,638}
\newcommand{\ModelGrowthPercent}{35}
\newcommand{\BodyTimesWordCap}{Sixteen}
\newcommand{\FurtherQuestionsWordCap}{Twelve}
\newcommand{\VerdictEight}{contradicted}
\newcommand{\VerdictSixtyFour}{mixed}
\newcommand{\HeldOne}{held}
\newcommand{\HeldTwo}{did not hold}
\newcommand{\HeldThree}{did not hold}
\newcommand{\HeldFour}{held}
\newcommand{\GapEqualSixtyFour}{9.9}
\newcommand{\BodyGainEqualSixtyFour}{7.7}
\newcommand{\GapUnequalSixtyFour}{23.0}
\newcommand{\BodyGainUnequalSixtyFour}{3.9}
\newcommand{\LargestGapUpToThirtyTwo}{3.0}
\newcommand{\GapEqualRepeatHigh}{12.6}
\newcommand{\GapEqualRepeatLow}{7.2}
\newcommand{\BodyGainEqualRepeatHigh}{10.6}
\newcommand{\BodyGainEqualRepeatLow}{4.7}
\newcommand{\LearnTimesEight}{11}
\newcommand{\LearnSmallEight}{200}
\newcommand{\LearnLargeEight}{2,263}
\newcommand{\LearnTimesEightUnequal}{8}
\newcommand{\LearnTimesSixteen}{16}
\newcommand{\LearnTimesSixteenUnequal}{16}
\newcommand{\LearnTimesThirtyTwo}{23}
\newcommand{\LearnTimesThirtyTwoUnequal}{45}
\newcommand{\TimeLines}{24}
\newcommand{\TimeSame}{15}
\newcommand{\TimeFaster}{6}
\newcommand{\TimeSlower}{3}
\newcommand{\FactsKnownEight}{200}
\newcommand{\QuestionsKnownEight}{200}
\newcommand{\NewKnownEight}{193}
\newcommand{\NotKnownEight}{200}
\newcommand{\FactsLowestEight}{100.0}
\newcommand{\FactsKnownSixteen}{120}
\newcommand{\QuestionsKnownSixteen}{120}
\newcommand{\NewKnownSixteen}{117}
\newcommand{\NotKnownSixteen}{118}
\newcommand{\FactsLowestSixteen}{100.0}
\newcommand{\FactsKnownThirtyTwo}{120}
\newcommand{\QuestionsKnownThirtyTwo}{106}
\newcommand{\NewKnownThirtyTwo}{70}
\newcommand{\NotKnownThirtyTwo}{66}
\newcommand{\FactsLowestThirtyTwo}{97.0}
\newcommand{\FactsKnownSixtyFour}{22}
\newcommand{\QuestionsKnownSixtyFour}{14}
\newcommand{\NewKnownSixtyFour}{1}
\newcommand{\NotKnownSixtyFour}{12}
\newcommand{\FactsLowestSixtyFour}{76.7}
\newcommand{\RunsSmallThirtyTwo}{30}
\newcommand{\RunsLargerThirtyTwo}{90}
\newcommand{\NotLargerThirtyTwo}{36}
\newcommand{\RunsLargerUnequalThirtyTwo}{45}
\newcommand{\NotLargerUnequalThirtyTwo}{3}
\newcommand{\NotNarrowUnequalThirtyTwo}{66}
\newcommand{\NotWideUnequalThirtyTwo}{87}
\newcommand{\RunsPerPointEight}{25}
\newcommand{\RunsPerPointLater}{15}
\newcommand{\RunsPerPointSixtyFour}{6}
\newcommand{\NewSeenLargestDifference}{1}
\newcommand{\NewSmallSixtyFourEqual}{91}
\newcommand{\SeenSmallSixtyFourEqual}{100}
\newcommand{\NewSmallSixtyFourUnequal}{89}
\newcommand{\SeenSmallSixtyFourUnequal}{100}
\newcommand{\FactsEqualMidNarrow}{100.0}
\newcommand{\FactsEqualMidMiddle}{100.0}
\newcommand{\FactsEqualMidWide}{100.0}
\newcommand{\FactsEqualLargeNarrow}{90.1}
\newcommand{\FactsEqualLargeMiddle}{94.6}
\newcommand{\FactsEqualLargeWide}{97.8}
\newcommand{\FactsUnequalMidNarrow}{89.2}
\newcommand{\FactsUnequalMidMiddle}{90.0}
\newcommand{\FactsUnequalMidWide}{94.7}
\newcommand{\FactsUnequalLargeNarrow}{77.0}
\newcommand{\FactsUnequalLargeMiddle}{78.4}
\newcommand{\FactsUnequalLargeWide}{80.9}
\newcommand{\SettingsSixtyFour}{18}
\newcommand{\SettingsBelow}{8}
\newcommand{\SettingsBelowWord}{eight}
\newcommand{\RunsBelowSixtyFour}{14}
\newcommand{\LateGainLow}{1.4}
\newcommand{\LateGainHigh}{3.6}
\newcommand{\RepeatSpreadHigh}{5.8}
\newcommand{\LowGroupsWord}{three}
\newcommand{\BodyStepsCounted}{6}
\newcommand{\BodyStepsLonger}{5}
\newcommand{\BodyStepsAll}{11}
\newcommand{\BodyRiseLow}{3.9}
\newcommand{\BodyRiseHigh}{7.7}
\newcommand{\BodyRiseLowWhole}{4}
\newcommand{\BodyRiseHighWhole}{8}
\newcommand{\LongerReached}{2}
\newcommand{\LongerReachedWord}{two}
\newcommand{\LongerShortWord}{six}
\newcommand{\ReachedBeforeLow}{94.6}
\newcommand{\ReachedBeforeHigh}{94.7}
\newcommand{\ShortEndLow}{83.3}
\newcommand{\ShortEndHigh}{93.8}
\newcommand{\LongerLowest}{83}
\newcommand{\LongerLowestExact}{83.3}
\newcommand{\LongerGainLow}{3.5}
\newcommand{\LongerGainHigh}{7.9}
\newcommand{\DoublingSecond}{289,631}
\newcommand{\DoublingThird}{579,262}
\newcommand{\DoublingEarlierLow}{2.0}
\newcommand{\DoublingEarlierHigh}{3.8}
\newcommand{\DoublingLaterLow}{1.2}
\newcommand{\DoublingLaterHigh}{3.7}
\newcommand{\LongerSlowing}{7}
\newcommand{\LastStretch}{60,738}
\newcommand{\LongerLevelWord}{two}
\newcommand{\LongerStillWord}{six}
\newcommand{\LongerStillWordCap}{Six}
\newcommand{\LastGainLow}{0.2}
\newcommand{\LastGainHigh}{0.5}
\newcommand{\LowestShort}{11.7}
\newcommand{\LowestPace}{2.3}
\newcommand{\LowestDoublingsNeeded}{five}
\newcommand{\LowestTimesNeeded}{32}
\newcommand{\LongerGapSmall}{16.7}
\newcommand{\LongerQuestionsLow}{66}
\newcommand{\LongerQuestionsHigh}{84}
\newcommand{\LongerNewLow}{60}
\newcommand{\LongerNewHigh}{79}
\newcommand{\LongerNotLow}{29}
\newcommand{\LongerNotHigh}{61}
\newcommand{\SharedDifferenceHigh}{3.0}
\newcommand{\SharedDifferenceOthers}{1.8}
\newcommand{\OtherLongerWord}{seven}
\newcommand{\LongerThreadsDifferWord}{six}
\newcommand{\ExampleRepeatLow}{87.4}
\newcommand{\ExampleRepeatHigh}{92.8}
\newcommand{\ExampleLongerShared}{84.1}
\newcommand{\ExampleLongerLate}{89.0}
\newcommand{\ExampleLongerEnd}{88.9}
\newcommand{\SixtyFourRows}{equal & 16 & 8 & 100.0 & 100.0 & 100.0 & 100 & 100.0 & \none & \none & \none & \none \\  &  & 32 & 100.0 & 100.0 & 100.0 & 100 & 100.0 & \none & \none & \none & \none \\  &  & 128 & 100.0 & 100.0 & 100.0 & 100 & 100.0 & \none & \none & \none & \none \\ \addlinespace[3pt]  & 256 & 8 & 100.0 & 100.0 & 100.0 & 84 & 100.0 & \none & \none & \none & \none \\  &  & 32 & 100.0 & 100.0 & 100.0 & 88 & 100.0 & \none & \none & \none & \none \\  &  & 128 & 100.0 & 100.0 & 100.0 & 99 & 100.0 & \none & \none & \none & \none \\ \addlinespace[3pt]  & 1,024 & 8 & 87.4 & 92.8 & 90.1 & 66 & 87.1 & 84.1 & 88.9 & +4.9 & 66 \\  &  & 32 & 91.7 & 97.5 & 94.6 & 70 & 91.4 & 91.0 & \reached{95.4} & +4.4 & 84 \\  &  & 128 & 98.0 & 97.5 & 97.8 & 81 & 97.7 & \none & \none & \none & \none \\ \midrule unequal & 16 & 8 & 100.0 & 100.0 & 100.0 & 100 & 100.0 & \none & \none & \none & \none \\  &  & 32 & 100.0 & 100.0 & 100.0 & 100 & 100.0 & \none & \none & \none & \none \\  &  & 128 & 100.0 & 100.0 & 100.0 & 100 & 100.0 & \none & \none & \none & \none \\ \addlinespace[3pt]  & 256 & 8 & 87.4 & 91.1 & 89.2 & 65 & 87.3 & 87.1 & 92.6 & +5.5 & 72 \\  &  & 32 & 88.9 & 91.1 & 90.0 & 68 & 88.1 & 89.7 & 93.8 & +4.1 & 73 \\  &  & 128 & 95.0 & 94.3 & 94.7 & 74 & 94.7 & 94.6 & \reached{98.0} & +3.5 & 79 \\ \addlinespace[3pt]  & 1,024 & 8 & 77.2 & 76.7 & 77.0 & 53 & 77.0 & 77.6 & 83.3 & +5.7 & 66 \\  &  & 32 & 79.3 & 77.6 & 78.4 & 54 & 78.5 & 79.4 & 85.8 & +6.4 & 73 \\  &  & 128 & 82.5 & 79.3 & 80.9 & 69 & 82.0 & 80.2 & 88.1 & +7.9 & 80 \\}


\begin{document}

% ------------------------------------------------------------------------------------------------
% Page one: the question, the answer, the picture
% ------------------------------------------------------------------------------------------------

\smallcaps{SECOND EXPERIMENT \enspace·\enspace REPORT OF 23 NOVEMBER 2026}\par
\vspace{4pt}
{\sffamily\bfseries\fontsize{21.5}{25}\selectfont Does a larger vocabulary need a larger model?\par}
\vspace{4pt}
{\sffamily\fontsize{10.5}{14}\selectfont\color{second}The claim under test is the owner's: a model of a given
size can hold the meaning of only so many words, so a larger vocabulary at the same model size does
worse.\par}
\vspace{5pt}{\color{hair}\hrule height 0.6pt}\vspace{6pt}

\lead{The test.} Small language models learned a made-up language about creatures and were then
tested on what they knew: \RunsCounted{} runs for the results, the \term{counted runs}, and
\RunsLonger{} longer ones. Four things were varied: the number of creatures, which is the size of
the vocabulary (16 to 1,024); the size of the model, its \term{body width} (8 to 128); how much
there is to know about each creature, in \term{bits} (8 to 64); and the \term{word use}: every
creature equally often (equal use), or a few far more often than the rest (unequal use). The main
score is \term{facts known}: 0 to 100, where a guess scores 25 and \KnownBar{} or more counts as
known. Section~1 shows all of this in pictures.

\begin{tcolorbox}[enhanced,colback=pale,colframe=pale,boxrule=0pt,arc=0pt,left=11pt,right=11pt,top=6pt,
  bottom=7pt,borderline west={2.2pt}{0pt}{ink},before skip=7pt,after skip=2pt]
\smallcaps{THE ANSWER}\par
\vspace{0.5pt}
\begin{enumerate}[leftmargin=1.35em,itemsep=2.5pt,topsep=1pt,label=\textbf{\arabic*.}]
\item \lead{Up to 32 bits per creature: no, for the facts.} The smallest model knew the facts of
      1,024 creatures as completely as those of 16: all \RunsEightSixteen{} runs at 8 and 16 bits
      ended at 100, and the lowest of the \RunsThirtyTwo{} runs at 32 bits at
      \FactsLowestThirtyTwo. The harder parts of the test began to slip at 32 bits, and more with
      more creatures. \see{findings 1 and 3}
\item \lead{More creatures cost training time instead.} At 8 bits, \CreatureTimes{} times the
      creatures needed about \LearnTimesEight{} times the training with equal use and
      \LearnTimesEightUnequal{} times with unequal use. \see{finding 2}
\item \lead{At 64 bits per creature: yes, in part.} 16 creatures were learned completely.
      With 1,024 creatures and equal use, the middle of two repeats was \FactsEqualLargeNarrow,
      \FactsEqualLargeMiddle{} and \FactsEqualLargeWide{} at body widths 8, 32 and 128: only the
      widest passed \KnownBar. With unequal use: \FactsUnequalLargeNarrow,
      \FactsUnequalLargeMiddle{} and \FactsUnequalLargeWide. \see{finding 4}
\item \lead{\LongerTimesWordCap{} times the training did not close that gap.} Each of the
      \SettingsBelow{} settings below \KnownBar{} was trained once more, as a new run
      \LongerTimesWord{} times as long. All \RunsLonger{} gained, ever more slowly. Only the
      \LongerReached{} that had been within half a point of \KnownBar{} reached it; the lowest
      ended at \LongerLowestExact. \see{finding 5}
\end{enumerate}
\end{tcolorbox}

\fighere{fig01_result_map}{fig:map}{\point{Every setting learned the facts, except where many creatures
meet 64 bits, and \LongerTimesWord{} times the training changed that little.} Each tile is one
setting: a number of creatures and a body width. A setting was trained \RepeatsFewest{} to
\RepeatsMost{} times, its repeats, and its tile shows the middle of their scores: the darker, the
more facts were missed at the end of training (a step is one round of training). The last panel
holds the \SettingsBelowWord{} settings whose counted tile is below \KnownBar. A longer run is not a
counted run carried on but a new run, and two runs of one setting can end several points apart. So
read its tile by itself: where that run stood at step \StepsShared, just short of the counted
length, and where it ended.}

% ------------------------------------------------------------------------------------------------
\section{How the experiment is built}
% ------------------------------------------------------------------------------------------------

\subsection{The language}

The language is about made-up creatures. Each creature has a fixed set of facts, and each fact is
one of four possible answers (Figure~\ref{fig:language}, left). Picking one answer out of four takes
two yes-or-no choices, so one fact is worth 2 \term{bits}, and a creature with four kinds of fact
carries 8 bits. To give a creature more to carry, more kinds of fact are added in the same pattern:
8, 16 or 32 kinds, which is 16, 32 or 64 bits per creature.

A model is never shown the facts. It only reads sentences about lists of one to four creatures:
statements, ``not'' statements, and questions with their answer (Figure~\ref{fig:language}, right).

\fig{fig02_language}{fig:language}{\point{A creature is a word with a set of facts, and a model only
ever reads sentences made from them.} Left: the facts of one creature in the language with 16
creatures. Right: the three kinds of sentence, and a pair of creatures kept apart for the test.}

\subsection{The test}

After training, each model sits a fixed test with four parts (Figure~\ref{fig:scoring}). All four
scores run from 0 to 100. Facts known is the main score, and the one the predictions are about. The
test asks every fact of every creature once, however rarely that creature appeared in training. In
this report a score of \KnownBar{} or more counts as ``known''.

\fig{fig03_scoring}{fig:scoring}{\point{The test asks for the same facts in four ways.} In each card
the model is given the words in the boxes and must supply the next one. A question score is the
share of ``yes'' questions answered right plus the share of ``no'' questions answered right, minus
100. So always giving the same answer scores 0, and \QuestionExampleRight{} right in 100 of each kind
scores \QuestionExampleScore.}

\subsection{The model}

The model is a small \term{transformer}, the kind of model behind today's language systems. It reads
a sentence word by word and guesses the next word. \term{Training} is this guessing, repeated: after
each round of guesses, the model's numbers are nudged so that the right words would have scored a
little higher. One such nudge is a \term{training step}; the model reads \SentencesPerStep{}
sentences for each.

The model has two parts (Figure~\ref{fig:model}). The \term{word table} holds a row of 32 numbers
for every word, and these rows are trained with everything else. The \term{body} is two layers that
turn the rows of the words read so far into a guess. Each layer has an attention part, which lets a
word gather information from the words before it, and a middle part made of \term{neurons}: small
units that each add up what comes in and pass the result on. Neither part replaces a word's 32
numbers: what it works out is added to them. The \term{body width} is the number of neurons in each
middle part. It is what ``model size'' means in this experiment: 8, 16, 32, 64 or 128. The first
experiment, which ended undecided, changed the length of every word's row and the body together.
Here a row has 32 numbers in every run, and only the body changes.

\fig{fig04_model}{fig:model}{\point{A model looks up a row of numbers for each word, and its body
turns those rows into a guess at the next word.} Every model starts from the same word table
(lower left). Body width is the waist of the shapes at the lower right: the number of neurons in a
middle part. What a middle part works out is added to a word's 32 numbers, which also carry on
unchanged. So a narrow body limits how much work a layer can do, not what can get past it.}

\subsection{What was varied, and which runs were made}

Four things were varied: the number of creatures (16, 64, 256, 1,024), the body width (8 to 128),
the bits per creature (8, 16, 32, 64), and the word use. With equal use every creature appears
equally often in the sentences. With unequal use the second most used creature appears half as
often as the first, the third a third as often, and so on, much as words do in real text.

One combination of the four is a \term{setting}. Each setting was trained several times: its
\term{repeats}. A repeat starts from other random numbers in the body and reads other sentences;
the word table starts the same every time. Where the report gives one score for a setting, it is
\term{the middle of its repeats}: half of them scored higher and half lower (with two repeats, the
point halfway between them).

The runs were made in \StagesWord{} stages, one for each number of bits (Figure~\ref{fig:runs}). A run at 64
bits takes far longer than one at 8 bits, so that stage was cut to a smaller grid to stay inside the
cloud credit. Each stage has its own training length, and the 64-bit stage also a lower
\term{learning rate} (how large a nudge each training step makes). Both were fixed from practice runs
before the stage's counted runs began (Appendix~A). Only counted runs enter the results.

The plan also fixed one check beforehand. A low score can mean ``cannot hold it'' or ``has not
learned it yet'', so every setting that ended below \KnownBar{} was to get one run
\LongerTimesWord{} times as long. That happened in \SettingsBelow{} settings, all at 64 bits: the
\RunsLonger{} longer runs. A longer run is not a counted run carried on: it is a new run that
starts from the same numbers and reads the same sentences as the first repeat of its setting, yet
does not follow it exactly (finding~5). Longer runs are always orange;
Figures~\ref{fig:predictions} to~\ref{fig:scores} show counted runs only.

\fig{fig05_runs_map}{fig:runs}{\point{The \RunsCounted{} counted runs in \StagesWord{} stages, and
the \RunsLonger{} longer runs at 64 bits.} One dot is one run and one box is one setting, with equal
use in its left column and unequal use in its right. An empty place at 64 bits is a setting cut
from that stage. A diamond is the longer run of a setting whose middle ended below \KnownBar. The
\PracticeRuns{} practice runs are not shown (Appendix~A).}

% ------------------------------------------------------------------------------------------------
\section{What the runs showed}
% ------------------------------------------------------------------------------------------------

\labelled{BEFORE THE RUNS}{Four predictions and a rule for the verdict}

Before the first counted run, the claim was turned into three predictions about facts known at the
end of training, P1 to P3. The plan set the opposite expectation beside them, as P4.

\begin{description}[leftmargin=2.0em,labelwidth=1.5em,labelsep=0.5em,font=\sffamily\bfseries,itemsep=1pt,topsep=2pt]
\item[P1] A small vocabulary needs only a small model: with 16 creatures every body width reaches \KnownBar.
\item[P2] A larger vocabulary at the same size does worse: at body width 8, 1,024 creatures score at
      least 10 points below 16 creatures.
\item[P3] A larger model holds more: on 1,024 creatures, body width 128 scores at least 10 points
      above body width 8.
\item[P4] With a word table that is trained, the number of creatures makes little difference: at
      every body width, 1,024 creatures score within 10 points of 16 creatures.
\end{description}

P2 and P4 cannot both hold. P4 added that if anything failed, it would follow the bits per creature
and the body width, not the number of creatures (finding~2 gives its reason).

The rule for the verdict was fixed too. A prediction holds when it comes out that way in at least 4
of the 5 repeats (repeat~1 of one setting is compared with repeat~1 of the other, and so on). The
claim is \emph{supported} if P1, P2 and P3 hold with equal and with unequal use,
\emph{contradicted} if P4 holds (the plan names no word use there; this report asks for both), and
\emph{mixed} otherwise.

\finding{At 8 bits the claim was contradicted: the vocabulary made no difference}

All \RunsEight{} runs at 8 bits ended at 100 facts known: with 16 creatures or with 1,024, with body
width 8 or 128, with equal or with unequal use. So P1 \HeldOne, P2 and P3 \HeldTwo, and P4
\HeldFour{} (Figure~\ref{fig:predictions}, top row of each panel). By the rule the claim is
\textbf{\VerdictEight} at 8 bits, with both word uses.

\fighere{fig06_predictions}{fig:predictions}{\point{The two predictions that set the claim apart, P2 and
P3, failed at 8, 16 and 32 bits; at 64 bits the gap that P2 predicts opens.} Counted runs only. A
mark is the middle of the repeats of one stage. The grey band is where the prediction says it would
fall. Ticks and crosses: the rule, which was written for the 8-bit stage.}

The rule was written for the 8-bit stage, the only one with \RepeatsEightWord{} repeats. The lower
rows of Figure~\ref{fig:predictions} apply the same four measures to the later stages, on the middle
of the repeats. At 16 and 32 bits nothing changes for the facts (the largest gap between 16 and
1,024 creatures is \LargestGapUpToThirtyTwo{} points). At 64 bits the gap opens (findings~4 and~5).

\finding{More creatures cost training time, not model size}

What the vocabulary did change is how long learning took (Figure~\ref{fig:curves}). At 8 bits with
equal use, the middle of all runs of a vocabulary first knew the facts (\KnownBar{} of 100) at step
\LearnSmallEight{} with 16 creatures and at step \LearnLargeEight{} with 1,024: \CreatureTimes{}
times the creatures took about \LearnTimesEight{} times the training. With more to know the same
step up cost more: about \LearnTimesSixteen{} times at 16 bits and \LearnTimesThirtyTwo{} times at
32 bits; with unequal use \LearnTimesEightUnequal, \LearnTimesSixteenUnequal{} and
\LearnTimesThirtyTwoUnequal{} times. Scores were measured at set steps, each \StepRatioShort{} times
the last, so these factors are rough.

\fighere{fig07_learning_curves}{fig:curves}{\point{A larger vocabulary is learned later, not worse; only
at 64 bits had the larger ones barely arrived, or not yet, when training ended.} Facts known during
training: each of the \RunsCounted{} counted runs as a thin line, the middle of the runs of each
vocabulary as a bold one (Figure~\ref{fig:longer} carries the 64-bit curves on). The 64-bit stage
had no 64-creature language and a lower learning rate, so its speed cannot be compared.}

The body width did not change these times in any steady way (Figure~\ref{fig:time}). Of the
\TimeLines{} lines up to 32 bits, the widest body took as long as the narrowest in \TimeSame, less
long in \TimeFaster{} and longer in \TimeSlower.

\fighere{fig08_learning_time}{fig:time}{\point{The time to learn rises with the number of creatures and
not with the body width: the lines lie one above the other and run about level.} The first measured
step at which the middle of a setting's repeats knew the facts; each mark is ten times the one
below. ``Not yet'': not at \KnownBar{} when training ended. At 64 bits only body widths 8, 32 and
128 were run, at a lower learning rate.}

\lead{A likely reason the vocabulary did not strain the model} is in the word table
(Figure~\ref{fig:model}, lower left). Every creature has a row of 32 trained numbers of its own, and
8 bits fit easily into 32 numbers. So a new word brings its own storage: at 8 bits the model for 16
creatures uses \RowsSmall{} rows of the table and the model for 1,024 creatures all \RowsLarge,
which hold \TableNumbersLarge{} numbers; the rest of the model has \OutsideNarrow{} at body width 8
and \OutsideWide{} at 128. What a new word still needs is to be met often enough, and that takes
training time. The runs do not prove this reason; they only fit it.

\finding{More bits per creature strain a model; the harder questions give way first}

Figure~\ref{fig:scores} puts the four parts of the test side by side. Up to 16 bits every score stays
at the top for every vocabulary. From 32 bits on, the abilities give way in a clear order.

\fighere{fig09_scores_by_bits}{fig:scores}{\point{The scores hold up to 16 bits, begin to slip at 32 and
fall at 64, the plain facts last and the smallest vocabulary last.} The middle of all counted runs
of each vocabulary: \RunsPerPointEight{} runs at 8 bits, \RunsPerPointLater{} at 16 and 32 bits,
\RunsPerPointSixtyFour{} at 64 bits (three body widths, no 64-creature language). Each stage had
its own training length, and the 64-bit stage its own learning rate, so a step between stages mixes
``more bits'' with ``trained differently''.}

\begin{itemize}
\item \lead{First the ``not'' sentences and the questions about new combinations of creatures.} At 32 bits
      \NotKnownThirtyTwo{} of \RunsThirtyTwo{} runs scored \KnownBar{} or more on ``not'' and
      \NewKnownThirtyTwo{} on new combinations; at 64 bits, \NotKnownSixtyFour{} and
      \NewKnownSixtyFour{} of \RunsSixtyFour.
\item \lead{Then the questions about one creature:} \QuestionsKnownThirtyTwo{} of \RunsThirtyTwo{}
      runs at 32 bits, \QuestionsKnownSixtyFour{} of \RunsSixtyFour{} at 64 bits.
\item \lead{The plain facts last.} All \RunsThirtyTwo{} runs at 32 bits still knew \KnownBar{} or more. At 64
      bits \FactsKnownSixtyFour{} of \RunsSixtyFour{} did.
\item \lead{And the larger vocabularies first.} At 32 bits all \RunsSmallThirtyTwo{} runs with 16
      creatures scored \KnownBar{} or more on ``not''; with 64 creatures or more,
      \NotLargerThirtyTwo{} of \RunsLargerThirtyTwo{} did, and with unequal use
      \NotLargerUnequalThirtyTwo{} of \RunsLargerUnequalThirtyTwo. A wider body helped there: with
      unequal use their middle ``not'' score was \NotNarrowUnequalThirtyTwo{} at body width 8 and
      \NotWideUnequalThirtyTwo{} at 128. So on the harder questions the claim's pattern begins to show
      at 32 bits, a stage before the facts.
\end{itemize}

\lead{It is not the newness of a combination that makes those questions hard.} The same kind of
question about combinations that did occur in training scored within \NewSeenLargestDifference{}
point, in the middle of all runs of a vocabulary from 64 creatures up: what is hard is a question
about several creatures at once. Only with 16 creatures at 64 bits did new combinations score lower
(\NewSmallSixtyFourEqual{} against
\SeenSmallSixtyFourEqual{} with equal use, \NewSmallSixtyFourUnequal{} against
\SeenSmallSixtyFourUnequal{} with unequal use).

\lead{A useful way to read the bottom axis of Figure~\ref{fig:scores}:} a creature's row has 32 numbers. At 32 bits each number
has to carry one bit, at 64 bits two. The scores begin to slip at about one bit per number. Whether
longer rows move the limit is what the third experiment asks (section~4).

\finding{At 64 bits the vocabulary matters, and a wider body helps by \BodyRiseLowWhole{} to \BodyRiseHighWhole{} points}

Here, for the first time, the number of creatures decides how many of the plain facts a model knows
at the end (Figure~\ref{fig:sixtyfour}). With 16 creatures every run ended at 100. With 256
creatures every setting ended at 100 with equal use, and at \FactsUnequalMidNarrow{} to
\FactsUnequalMidWide{} with unequal use. With 1,024 creatures the settings ended at
\FactsEqualLargeNarrow{} to \FactsEqualLargeWide{} with equal use and at \FactsUnequalLargeNarrow{}
to \FactsUnequalLargeWide{} with unequal use.

\fighere{fig10_sixty_four}{fig:sixtyfour}{\point{At 64 bits more creatures mean fewer facts known, a wider
body helps by a few points, and the longer runs gained a few points without closing the gap.} All
\SettingsSixtyFour{} settings. The two dots are the two counted runs of a setting, side by side
where one would cover the other; the thin line joins the middles. Where a middle is below
\KnownBar{} the setting was trained once more: the orange arrow runs from where that new run stood
at step \StepsShared, its last measurement before the counted length, to where it ended at step
\StepsLonger, with its final score. The foot of an arrow is the new run's own score, not a counted
one. At that step it lay within \SharedDifferenceOthers{} points of the setting's first repeat in
\OtherLongerWord{} settings and \SharedDifferenceHigh{} points below it in one: 1,024 creatures,
equal use, body width 8 (Appendix~C). The axis starts at \AxisStartSettings.}

\begin{itemize}
\item \lead{A wider body helped, by a few points.} In the \LowGroupsWord{} groups that fell short (1,024
      creatures with equal use, 256 and 1,024 with unequal use), every step up in body width raised
      the score, by \BodyRiseLow{} to \BodyRiseHigh{} points from width 8 to 128. Two repeats of a
      setting ended up to \RepeatSpreadHigh{} points apart, so the size is uncertain; the direction
      is not, for the longer runs repeat it (finding~5).
\item \lead{Unequal use made it worse.} The test asks every fact of every creature once, so the
      least used of 1,024 creatures counts as much as the most used, which appears a thousand
      times as often. A likely reason for the lower scores, which these runs do not test: a rarely
      used creature is met far less often during training.
\item \lead{The pattern of the claim appears in part.} With unequal use, 1,024 creatures score
      \GapUnequalSixtyFour{} points below 16 creatures at body width 8, so P2 would hold. With equal
      use the stage falls on the 10 points that decide between P2 and P4. The gap is
      \GapEqualSixtyFour, the halfway point of two repeats that give \GapEqualRepeatHigh{} and
      \GapEqualRepeatLow: one supports P2 and the other P4. And the facts of 1,024 creatures were
      known only at body width 128, in both repeats. That is the claim's pattern in this one row,
      though the gain over width 8, \BodyGainEqualSixtyFour{} points (\BodyGainEqualRepeatHigh{}
      and \BodyGainEqualRepeatLow{} in the two repeats), is short of the 10 that P3 asks for; with
      unequal use it is \BodyGainUnequalSixtyFour. Read by the rule, which was fixed for \RepeatsEightWord{} repeats where this stage has
      \RepeatsSixtyFourWord, the stage would be \emph{\VerdictSixtyFour}.
\item \lead{One expectation written down beforehand was wrong here.} P4 added that failures would
      follow the bits and the body width, not the number of creatures. At 64 bits the number of
      creatures plainly matters.
\end{itemize}

\finding{\LongerTimesWordCap{} times the training: all \RunsLongerWord{} gained, only the \LongerReachedWord{} nearest to \KnownBar{} reached it}

When the counted 64-bit runs ended, the \RunsBelowSixtyFour{} runs below \KnownBar{} were all still rising:
each gained \LateGainLow{} to \LateGainHigh{} points between step \StepsLate{} and the end at
\StepsSixtyFour. That is the case the plan had provided for. Each of the \SettingsBelow{} settings
whose middle was below \KnownBar{} was run once more, for \StepsLonger{} steps (Figure~\ref{fig:longer}).

\fighere{fig11_longer_curves}{fig:longer}{\point{With \LongerTimesWord{} times the training every one of
the \RunsLongerWord{} gained, more slowly as it went on, and \LongerShortWord{} were still below
\KnownBar{} at the end.} Facts known during training at 64 bits. Unlike
Figure~\ref{fig:curves}, the steps are evenly spaced: the early learning is pressed against the
left edge, and the slowing afterwards shows. Blue: all \RunsSixtyFour{} counted runs, which end at
\StepsSixtyFour{} steps. Orange: the \RunsLonger{} longer runs, with a mark at each later
measurement and the final score at the end. A dashed line is unequal use, not a forecast. The axis
starts at \AxisStartCurves; every run begins near 25.}

\lead{How to read a longer run against the counted runs.} A longer run starts from the same numbers
and reads the same sentences as the first repeat of its setting. But it was made on another
computer, and \LongerThreadsDifferWord{} of the \RunsLongerWord{} with another number of processor
threads; either can change the last digits of a calculation. That is the likely reason, not tested
here, why the two runs begin alike and then drift apart, by up to
\SharedDifferenceHigh{} points at step \StepsShared. Take 1,024 creatures, equal use, body width 8.
Its two counted repeats ended at \ExampleRepeatLow{} and \ExampleRepeatHigh. Its longer run stood at
\ExampleLongerShared{} at step \StepsShared{} and ended at \ExampleLongerEnd: after
\LongerTimesWord{} times the training, still between the two counted repeats. So a longer run is
measured against its own earlier score, and Figure~\ref{fig:sixtyfour} draws it as an arrow of
its own.

\begin{itemize}
\item \lead{Every one gained.} Between step \StepsShared{} and the end, facts known rose by
      \LongerGainLow{} to \LongerGainHigh{} points. So part of what looked like a limit was learning
      that had not finished.
\item \lead{The gains shrink.} A doubling of the training, from step \StepsShared{} to step
      \DoublingSecond, gained \DoublingEarlierLow{} to \DoublingEarlierHigh{} points. The next
      doubling, to step \DoublingThird, gained \DoublingLaterLow{} to \DoublingLaterHigh; in
      \LongerSlowing{} of the \RunsLonger{} runs it brought less than the first. Over the last
      \LastStretch{} steps \LongerStillWord{} runs still gained \LastGainLow{} to \LastGainHigh{}
      points and \LongerLevelWord{} were level: 256 creatures with unequal use at body width 32,
      and 1,024 creatures with equal use at body width 8 (\ExampleLongerLate, then
      \ExampleLongerEnd).
\item \lead{Only the \LongerReachedWord{} nearest to \KnownBar{} reached it.} In those
      \LongerReachedWord{} settings the counted middle had been within half a point of it
      (\ReachedBeforeLow{} and \ReachedBeforeHigh). The other \LongerShortWord{} ended short, at
      \ShortEndLow{} to \ShortEndHigh. The lowest (1,024 creatures, unequal use, body width 8) was
      still \LowestShort{} points short: at the \LowestPace{} points of its last doubling, more
      than \LowestDoublingsNeeded{} further doublings, or over \LowestTimesNeeded{} times the
      \StepsLonger{} steps.
\item \lead{The order did not change.} After \LongerTimesWord{} times the training a wider body
      still scored higher in each of those \LowGroupsWord{} groups, 1,024 creatures still scored
      below 256 at every body width, and unequal use below equal use. With unequal use at body
      width 8, 1,024 creatures were still \LongerGapSmall{} points below 16 creatures. Each of these
      comparisons rests on single runs.
\item \lead{The harder questions stayed far from \KnownBar.} Questions about one creature ended
      between \LongerQuestionsLow{} and \LongerQuestionsHigh, questions about new combinations
      between \LongerNewLow{} and \LongerNewHigh, and ``not'' known between \LongerNotLow{} and
      \LongerNotHigh, where a guess scores 25.
\end{itemize}

% ------------------------------------------------------------------------------------------------
\section{What this shows, and what it does not}
% ------------------------------------------------------------------------------------------------

\lead{It shows} that in this kind of model, where every word has a trained row of its own, adding
words does not by itself fill the model up. Adding to what each word must carry does. And once the
words carry a lot, the number of words begins to count as well, first on the harder questions (at
32 bits) and then on the plain facts (at 64 bits), and the size of the model begins to help.

\lead{It does not show:}
\begin{itemize}
\item \lead{Whether 64 bits is a hard limit.} With 1,024 creatures the models had not learned
      everything after \LongerTimesWord{} times the training. \LongerStillWordCap{} of the
      \RunsLongerWord{} longer runs were still gaining, slowly, when they stopped, and
      \LongerLevelWord{} were level. Whether more training would ever get them there is open; the
      longer runs show only that it is a very slow way to close the gap.
\item \lead{Which part of the model sets the limit, or what a model larger in another way would
      do.} ``Model size'' was the body width only, and a word's row had 32 numbers in every run.
      \BodyTimesWordCap{} times the neurons make the whole model only \ModelGrowthPercent{}
      percent larger (from \ModelNarrow{} numbers to \ModelWide, with 1,024 creatures at 8 bits),
      because the word table is the larger part and stays as it is. Whether longer rows move the limit is
      the question of the third experiment; more layers were not tried.
\item \lead{Why the vocabulary did not matter at fewer bits.} ``Each word brings its own storage''
      fits the runs, but no run tested it. The test would be a fixed word table, in which only the
      body is trained.
\item \lead{A clean comparison across bits.} Each stage had its own training length, and the 64-bit
      stage its own learning rate (Appendix~A). A difference between two stages mixes ``more bits''
      with ``trained differently''.
\end{itemize}

\lead{How firm the 64-bit results are.} That stage has \RepeatsSixtyFourWord{} repeats of each
setting and one longer run for each of the \SettingsBelowWord{} low settings, and two repeats differed
by up to \RepeatSpreadHigh{} points. The broad pattern is far larger than that: 16 creatures at 100,
1,024 creatures well below, unequal use lowest. The direction of the body width's effect is firm
too: all \BodyStepsAll{} steps from one body width to the next go up, \BodyStepsCounted{} in the
counted runs and \BodyStepsLonger{} in the longer runs. Its size, \BodyRiseLowWhole{} to
\BodyRiseHighWhole{} points, is not firm, and neither is any difference of a few points between a
longer run and a counted run.

% ------------------------------------------------------------------------------------------------
\section{What comes next}
% ------------------------------------------------------------------------------------------------

\begin{itemize}
\item \lead{A third experiment is running now,} on rented cloud graphics cards (plan:
      \s{v3/PLAN.md}). It varies the length of a word's row (8, 16, 32 or 64 numbers), with body
      widths 4 to 128 and with 8 to 128 bits per creature: \ThirdRuns{} runs, three repeats of each
      setting. Every run has the same learning rate and ends by one fixed rule when it stops
      gaining, which removes this experiment's main weakness; every model of a language reads the
      same text. This report contains no results from it.
\item \lead{The run with a fixed word table} (section~3) is prepared in the programs and has not
      been made.
\item \lead{\FurtherQuestionsWordCap{} further questions} raised along the way are kept in
      \s{research\_questions.md}, among them whether the vocabulary matters when words depend on
      each other, whether variety in the text helps, and whether rewards can teach what plain
      training did not.
\end{itemize}

% ------------------------------------------------------------------------------------------------
\appendix
% The appendix follows on the same page, set off by a rule and its label.
\vspace{16pt plus 4pt minus 2pt}{\color{hair}\hrule height 0.6pt}\vspace{8pt}
\smallcaps{APPENDIX}\par\vspace{-9pt}
\section{Settings, and how they were fixed}
% ------------------------------------------------------------------------------------------------

Each stage's learning rate and training length were fixed from practice runs before its counted
runs started. Practice runs use a starting point of their own and are never counted. All used equal
use; those at 16, 32 and 64 bits used 1,024 creatures, with body widths 8 and 128.

\begin{center}\sffamily\footnotesize
\begin{tabular*}{\linewidth}{@{\extracolsep{\fill}}lrrrrrl@{}}
\toprule
 & & & Training & Sentences read, & Learning & \\
 & Runs & Repeats & steps & in millions & rate & Practice runs behind the choice \\
\midrule
8 bits & \RunsEight & \RepeatsEight & \StepsEight & \SentencesEight & \RateEight & \PracticeEight: five settings, three rates \\
16 bits & \RunsSixteen & \RepeatsSixteen & \StepsSixteen & \SentencesSixteen & \RateSixteen & \PracticeSixteen: two body widths, two rates \\
32 bits & \RunsThirtyTwo & \RepeatsThirtyTwo & \StepsThirtyTwo & \SentencesThirtyTwo & \RateThirtyTwo & \PracticeThirtyTwo: two body widths, two rates \\
64 bits & \RunsSixtyFour & \RepeatsSixtyFour & \StepsSixtyFour & \SentencesSixtyFour & \RateSixtyFour & \PracticeSixtyFour: two body widths, three rates \\
\addlinespace[2pt]
64 bits, longer runs & \RunsLonger & 1 & \StepsLonger & \SentencesLonger & \RateSixtyFour & none: \LongerTimesWord{} times the counted length \\
\bottomrule
\end{tabular*}
\end{center}

\begin{itemize}
\item \lead{Learning rate.} The one with the most facts known at the end of practice; questions
      broke a tie.
\item \lead{Training length.} Several times what the slowest practice run needed. At 64 bits no
      practice run had finished learning, so the length was set at four times the practice length.
\item \lead{The same in every run:} two layers, two attention heads (each attention part is
      split in two that look back separately), 32 numbers a word, \SentencesPerStep{} sentences a
      training step, one sentence a line.
\item \lead{When scores were measured.} At steps \FirstSteps{} and so on, each about \StepRatio{}
      times the last, and at the last step. So a counted 64-bit run was last measured at steps
      \CountedLastSteps, and a longer run at steps \LongerLastSteps: it has none at
      \StepsSixtyFour, so the report compares the two at step \StepsShared.
\item \lead{Sizes.} Outside the word table a model has \OutsideNarrow{} trained numbers at body width
      8 and \OutsideWide{} at body width 128 (at 8 bits). The word table adds 32 for every word in
      use. At 8 bits it has \TableRowsEight{} rows: the 1,024 creature names, \OtherWords{} other
      words and marks, and one filler row, which pads short sentences. Every further kind of fact
      adds \WordsPerKind{} words, which gives \TableRowsSixteen, \TableRowsThirtyTwo{} and
      \TableRowsSixtyFour{} rows at 16, 32 and 64 bits. \RowsNeverChangedWordCap{} rows in use
      never moved in any run: the filler row, the full stop, \s{yes} and \s{no}. They only ever
      end a sentence, so nothing is guessed after them.
\end{itemize}

\section{Where it ran, what it cost, and what differed from the plan}

\begin{itemize}
\item \lead{Where.} On Modal, a cloud service, on up to \CloudMachines{} small machines at once, each
      with one or two processor cores and no graphics card. Of the 16-bit and 32-bit stages,
      \RunsOnOwnComputer{} runs were made on the owner's own computer with the same programs, and so
      were all \RunsLonger{} longer runs. A longer run took between \LongerHoursLow{} and
      \LongerHoursHigh{} hours.
\item \lead{Cost.} \$\CloudCost{} of cloud credit, by the service's own meter.
\item \lead{Different from the plan:}
  \begin{itemize}
  \item The 64-bit stage was cut from \PlannedSixtyFour{} runs to \RunsSixtyFour{} to stay inside the
        credit.
  \item Part of the 16-bit and 32-bit stages ran on the owner's computer, for the same reason.
  \item One word of the language was changed: the value \s{mice} was also a creature's name and
        became \s{gnats}.
  \item The longer runs were started on 22~November when \LongerStartRunsIn{} of the
        \RunsSixtyFour{} runs of the 64-bit stage were in; the last two (256 creatures, equal use,
        body width 128) ended at 100 and so changed nothing. The longer runs were stopped once that
        night and went on from their saved state, which the programs are built to do without
        changing the result. They ended on 23~November.
  \item \TimingRuns{} short timing runs and one tiny end-to-end test were made. They are not
        counted and not shown.
  \end{itemize}
\item \lead{The word table.} The first design had a fixed word table. The owner changed it to a
      trained one before any run. Both expectations were then written into the plan side by side,
      as P2 and P4.
\item \lead{Every counted run is reported.} Nothing was rerun to get a better number.
\end{itemize}

\section{The 64-bit settings in exact numbers}

Facts known for all \SettingsSixtyFour{} settings at 64 bits, and the score on questions about one
creature. A dash: not trained longer, because the middle of the counted runs was \KnownBar{} or
more. Bold: reached \KnownBar. A longer run starts from the same numbers as repeat 1 of its setting;
the two columns at step \StepsShared{} show how far it had drifted from that repeat, and its gain
is counted from its own score at that step. Gains are worked out before rounding, so one can differ
by 0.1 from what the printed scores give.

\begin{center}\sffamily\footnotesize
\setlength{\tabcolsep}{0pt}
% Twelve columns of fixed width; \gap is the space between the four groups.
\newcommand{\gap}{\hspace{0.24in}}
\renewcommand{\arraystretch}{1.12}  % a little air between the rows of a dense table
\begin{tabular}{@{}w{l}{0.56in}w{l}{0.60in}w{r}{0.32in}@{\gap}w{r}{0.48in}w{r}{0.48in}w{r}{0.42in}w{r}{0.52in}%
  @{\gap}w{r}{0.45in}w{r}{0.50in}@{\gap}w{r}{0.42in}w{r}{0.38in}w{r}{0.55in}@{}}
\toprule
 & & & \multicolumn{4}{@{\gap}l}{Counted runs at their end} & \multicolumn{2}{@{\gap}l}{At step \StepsShared}
   & \multicolumn{3}{@{\gap}l@{}}{Longer run at its end} \\
 & & & \multicolumn{4}{@{\gap}l}{(\StepsSixtyFour{} steps)} & \multicolumn{2}{@{\gap}l}{(facts known)}
   & \multicolumn{3}{@{\gap}l@{}}{(\StepsLonger{} steps)} \\
\cmidrule(l{0.24in}){4-7}\cmidrule(l{0.24in}){8-9}\cmidrule(l{0.24in}){10-12}
 & & Body & \multicolumn{3}{@{\gap}l}{facts known} & questions & & longer & facts & & \\
Word use & Creatures & width & repeat 1 & repeat 2 & middle & middle & repeat 1 & run & known & gain & questions \\
\midrule
\SixtyFourRows
\bottomrule
\end{tabular}
\end{center}

\end{document}
